\chapter{Conclusion and Future Work}
\label{ch:conclusion}

\section{Summary of contributions}
This thesis asked whether an agent can learn how to construct its own memory and
whether the learned optimum is task-dependent. The answer, established across
three stages of increasing realism, is yes. We parameterized memory construction
--- storage, abstraction, and retrieval scoring --- by a single vector
$\thetavec$ and optimized it by black-box search, leaving the policy and the
language model untouched. In grid worlds the optimizer recovered visibly
different $\thetavec$ per task and lifted the hardest task from $2.5\%$ to
$27.5\%$ success. On a twelve-system benchmark a ten-dimensional learnable memory
reached the top performance cluster, and a tuned $\thetavec$ generalized within
but not across task families. On six real long-context LLM benchmarks,
corpus-cumulative tuning produced a recency-independent memory whose
end-of-corpus accuracy lift survived held-out testing on two of the three
domain-coherent benchmarks (FinanceBench, CUAD; QASPER did not survive Holm
correction) and held when CUAD was scaled from ten to fifty contracts. The
tuning objective is a valid proxy for answer quality: recall@8 and judge score
correlate at pooled Spearman $\rho=0.69$.

The work's distinctive feature is that it audits itself. The headline that a
full-context baseline ``collapses'' was traced to a bug and corrected to an
accuracy \emph{tie}; the contribution of learned, task-adaptive memory at
this scale is therefore reframed as one of \textbf{efficiency and scalability}
--- matching full-context accuracy at roughly $1/18$th the cost --- rather than
raw accuracy. Concretely, on FinanceBench the dump-all and selective answers are
statistically non-inferior (overlapping bootstrap CIs, $0.607$--$0.689$ vs.\ the
tuned cell) while the selective prompt stays flat at $\sim$704 tokens where
dump-all exceeds the 128K context window at $N{\approx}11$ CUAD contracts and
reaches $43\times$ the tokens at the full $N{=}510$. The one-time CMA-ES tuning
cost --- no LLM in the loop, only \recallk{} --- was approximately \$18.5 of
compute per benchmark, amortized over all subsequent queries.

\section{Limitations}
The Stage-1 results are single-seed. Several Stage-3 cells are single-seed
(HotpotQA and LongMemEval were re-run at 100-document scale, but multi-seed
replication of the larger benchmarks remains future work; corpus-mode evaluation
is near-deterministic at temperature~0, so the gain would be CI tightening, not a
point-estimate change). CUAD is now evaluated on 50 of its 510 contracts for the
headline pair; the auxiliary baselines remain at the 10-contract pilot scale.
Answers are judged by a single (cross-vendor) Claude judge; an independent blind
second pass over a stratified 180-question sample gives quadratic-weighted
Cohen's $\kappa=0.66$ (\emph{substantial}), but because both passes are
Claude-class this bounds judge \emph{self-consistency}, not cross-vendor or human
agreement --- a fully independent second-vendor or human rater remains future
work. Refusal/acknowledgment answers carry validated rule-assisted scores
(population-weighted error $0.028$ on a 300-entry hand-checked sample) rather
than one-by-one judgments. Finally, all Stage-3 results use a single answerer
(\texttt{gpt-4o-mini}) and a single encoder (MiniLM); generalization across
models and encoders is untested.

\section{Future work}
Natural extensions are: (i) multi-seed replication of the larger Stage-3 cells;
(ii) scaling CUAD and the other benchmarks to their full corpora and all
configurations; (iii) a second independent judge with Cohen's $\kappa$; (iv)
revisiting the neural memory controller with a larger optimization budget, now
that the scalar \Vfour{} sets a strong, cheap baseline; (v) larger answerer
models, to test whether the efficiency advantage of learned selective memory
persists when the answerer is stronger; and (vi) making $\thetavec$ itself
adaptive \emph{within} an episode (a per-step controller) rather than fixed per
task, which the present grid results suggest could help on long-horizon
out-of-distribution tasks where a fixed $\thetavec$ fails.

\section{Closing}
Memory for LLM agents need not be hand-designed and frozen. A small parameter
vector that governs what to store and how to retrieve can be \emph{learned}, its
optimum is \emph{task-dependent}, and --- measured carefully --- it buys
efficiency and graceful scaling rather than a free accuracy gain. That is a
modest but defensible step toward agents that adapt not just their actions but
the structure of their own memory.
